\documentclass[12pt,a4paper]{article}
\usepackage[UTF8,heading=true]{ctex}
\usepackage{geometry}
\usepackage{amsmath,amssymb,amsthm}
\usepackage{graphicx}
\usepackage{booktabs}
\usepackage{longtable}
\usepackage{array}
\usepackage{enumitem}
\usepackage{hyperref}
\usepackage{fancyhdr}
\usepackage{xcolor}
\usepackage{listings}
\usepackage{setspace}  % 添加行距控制包

% 配置 ctex 章节格式为 "一、"
\ctexset{
  section = {
    name = {,、},
    number = \chinese{section},
    format = {\Large\bfseries\raggedright}
  }
}

\geometry{left=2.5cm,right=2.5cm,top=2.5cm,bottom=2.5cm}
\setlength{\headheight}{15pt}
\setstretch{1.5}  % 设置1.5倍行距（CNIPA推荐）

\pagestyle{fancy}
\fancyhf{}
\fancyhead[C]{发明专利技术交底书}
\fancyfoot[C]{\thepage}

\title{\textbf{发明专利技术交底书} \\[0.5em]
\Large 基于声子本征模复用的量子低密度奇偶校验码综合征存储与读取方法}
\author{中国移动通信有限公司研究院}
\date{\today}

\begin{document}

\begin{center}
    {\Large\heiti 中国移动专利申请} \\[1em]
    {\Huge\heiti 技术交底书} \\[1em]
\end{center}

\noindent
\begin{tabular}{|p{3cm}|p{12cm}|}
\hline
公司编号 &  \\
\hline
发明名称 & 基于声子本征模复用的量子低密度奇偶校验码综合征存储与读取方法 \\
\hline
申报单位 & 研究院 \\
\hline
申报类型 & 发明 \\
\hline
发明人 & 许达 \\
\hline
技术联系人 & 许达 (xudayj@chinamobile.com, +86-13521894156) \\
\hline
注意事项 & 1．技术联系人应为深入了解本申请提案技术方案的技术人员。 \\
         & 2．请按照集团公司提供的本技术交底书模板逐项填写。 \\
\hline
\end{tabular}

\vspace{1cm}

%==============================================================================
\section{发明名称}
%==============================================================================

基于声子本征模复用的量子低密度奇偶校验码综合征存储与读取方法


%==============================================================================
\section{技术领域}
%==============================================================================

本发明涉及量子计算和量子纠错技术领域，特别涉及一种面向线性离子阱量子计算机的 qLDPC 码综合征提取方法、编译方法及系统。

\textbf{与量子计算的关联说明：}

容错量子计算需要持续执行量子纠错。量子低密度奇偶校验码（qLDPC）能够以较高编码率存储大量逻辑量子比特，但其 Tanner 图校验节点通常对应跨越多个物理量子比特的高权重稳定子。在线性离子链中，传统做法需要为稳定子测量配置辅助离子承载奇偶信息，并通过多次 Mølmer-Sørensen 门或等效纠缠门顺序提取综合征。

本发明利用离子晶体的正交振动本征模作为临时综合征存储空间，将 qLDPC Tanner 图中的校验节点映射到可寻址声子模式，使多个重叠稳定子在声学子空间中并行写入，从而减少“每校验一个辅助离子”的空间开销并降低综合征提取深度。

\textbf{关联项目：}\underline{\hspace{8cm}}（待填写）

%==============================================================================
\section{现有技术的技术方案}
%==============================================================================

\subsection{离子阱量子纠错技术现状}

现有离子阱量子计算通常以内态能级表示量子比特，以集体振动模式作为实现多量子比特门的媒介。对于稳定子码，综合征提取一般通过辅助离子、控制 Pauli 门、MS 型纠缠门和辅助离子测量完成。该方案适用于小规模表面码、颜色码或低权重校验，但在 qLDPC 码中会遇到以下限制：

\begin{itemize}
    \item 高权重稳定子需要多个门层或较长全局脉冲序列；
    \item 重叠校验节点需要多个辅助离子或复杂调度，空间开销较大；
    \item 多个校验顺序执行时，数据量子比特在等待过程中继续退相干。
\end{itemize}

\subsection{声子模式利用现状}

线性 $N$ 离子链具有 $3N$ 个正交运动本征模，包括轴向和径向的质心、呼吸、摇摆或 zig-zag 类模式。受模式频率间隔、冷却效率、Lamb-Dicke 参数和激光谱泄漏约束，实际可并行使用的是其中的可寻址模式子集。现有技术主要将这些模式作为纠缠门的中介、量子模拟自由度或少量运动态编码资源，并未将 qLDPC Tanner 图中的多个校验节点系统性映射为多个可并行写入、可逐模读取的综合征存储单元。

%==============================================================================
\section{现有技术的缺点及本申请提案要解决的技术问题}
%==============================================================================

现有技术主要存在以下缺陷和技术问题：

(1) \textbf{辅助离子开销大}：若每个 qLDPC 校验节点使用一个专用辅助离子，则需要大量额外离子，并会增加装载、冷却、读出和串扰管理成本。

(2) \textbf{高权重校验执行慢}：权重为 6 至 12 的稳定子若按二体门或小规模 MS 门分解，电路深度高，退相干和运动加热误差累积明显。

(3) \textbf{重叠稳定子调度困难}：qLDPC Tanner 图中同一数据量子比特常参与多个校验，顺序测量会造成等待时间和误差传播，完全并行测量又容易发生辅助离子和激光寻址冲突。

(4) \textbf{运动模资源未被充分利用}：离子链天然具备多个正交本征模，但现有综合征提取一般只将其作为门媒介，而非作为独立的奇偶信息暂存寄存器。

本申请提案旨在解决上述问题，提供一种用声子本征模替代“每校验专用辅助离子”的 qLDPC 综合征提取方法。

%==============================================================================
\section{本申请提案的技术方案的详细阐述}
%==============================================================================

本发明提供一种 qLDPC 综合征提取编译器和执行方法，将待测稳定子集合 $\{S_j\}$ 映射到离子链可寻址正交本征模集合 $\{m_j\}$，并通过多色激光产生状态依赖力，将每个稳定子的奇偶值写入相应声子模的二值指针态或二值占据窗口。

\subsection{系统架构}

本发明的系统架构包括：

\textbf{离子链数据寄存器}：由 $N$ 个离子构成，承载 qLDPC 码的数据量子比特，并提供从 $3N$ 个运动本征模中筛选出的可用模式子集。

\textbf{qLDPC 编译模块}：读取 Tanner 图和稳定子矩阵，根据校验权重、模式频率间隔、Lamb-Dicke 参数和激光可寻址性，将校验节点分配至互相正交且频率可分辨的声子模式。

\textbf{多色状态依赖力控制模块}：对每个被选稳定子生成一组红/蓝边带或失谐驱动脉冲，使目标数据量子比特的联合奇偶值控制相应模式的条件位移、非目标模式位移闭合和二值化读出窗口。

\textbf{共享读出离子}：在离子链边缘设置一个或少量 Ba$^+$、Mg$^+$ 或其他双种类读出离子，通过红边带、蓝边带或量子逻辑读出脉冲逐个读取目标模式的二值指针态。

请参考\textbf{图 1}，其展示了系统整体架构及综合征写入、读取与译码反馈流程。
\begin{figure}[h]
    \centering
    \includegraphics[width=0.8\textwidth]{Figure_1.png}
    \caption{系统架构示意图}
    \label{fig:arch}
\end{figure}

\subsection{核心方法步骤}

\subsubsection{1. 稳定子到声子模式的编译映射}
对 qLDPC 校验矩阵 $H$ 中每个校验行 $h_j$，确定其支撑集合 $\mathrm{supp}(h_j)$。对于 CSS 码，分别调度 $X$ 型和 $Z$ 型校验；对于含 $Y$ 或混合 Pauli 项的稳定子，先对相关数据离子执行本地基旋转，使待测算符映射到统一测量基，再在读出后恢复或更新 Pauli frame。编译器从可用模式子集中选择模式 $m_j$，并约束任意两个同时执行的校验满足模式正交、频率间隔大于预设阈值、脉冲谱泄漏低于阈值且热占据低于初始化阈值；同时要求被选模式 $m_j$ 在校验支撑集合 $\mathrm{supp}(h_j)$ 内所有数据离子上均具有非可忽略且量级相近的参与系数，以保证奇偶可写入同一模式。无法同时满足约束的校验通过图着色方式分批调度。

如图 2 所示为 Tanner 图校验节点到声子模式的编译映射与分批调度示意。
\begin{figure}[h]
    \centering
    \includegraphics[width=0.8\textwidth]{Figure_2.png}
    \caption{Tanner 图校验节点到声子模式的编译映射示意图}
    \label{fig:mapping}
\end{figure}

\subsubsection{2. 并行声子写入}
对稳定子 $S_j=\prod_{i\in \mathrm{supp}(h_j)} P_i$，先通过本地基旋转将其归一化为 $\prod_i Z_i$ 形式。由于单段双色状态依赖力产生的耦合与支撑集合内各离子 $Z_i$ 算符的加权和（按模式参与向量加权）成正比，而非稳定子乘积本身，编译器将权重为 $w$ 的奇偶条件写入分解为若干段条件位移与中间载波旋转构成的复合脉冲序列。该复合序列的作用机制为：在相邻条件位移段之间插入作用于支撑集合内数据离子的 $\pi/2$ 载波旋转，将各离子单体 $Z_i$ 相位逐级折叠并干涉合成为支撑集合的联合奇偶可观测量（等效于以 Ramsey 型干涉或受控相位级联把 $Z_i$ 加权和折叠为 $\prod_i Z_i$），从而使位移幅度由稳定子乘积本征值而非单体本征值之和决定。各段脉冲在相空间中对非目标模式形成闭合回路，对目标模式 $m_j$ 则通过几何相位合成，使其末态仅由稳定子本征值（联合奇偶）决定。为保证奇偶能够写入同一模式，编译器要求支撑集合 $\mathrm{supp}(h_j)$ 内所有数据离子在目标模式 $m_j$ 上均具有非可忽略且量级相近的参与系数。该复合序列所合成的理想等效目标演化（而非单段双色脉冲的裸哈密顿量本身）可表示为
$$
U=\prod_j \exp\left[-i\theta_j \frac{1-S_j}{2}(a_j+a_j^\dagger)\right],
$$
其中 $a_j^\dagger$ 为模式 $m_j$ 的产生算符，$\theta_j$ 为由稳定子本征值控制的位移幅度，投影算符 $(1-S_j)/2$ 将系统投影到 $S_j=-1$ 的本征子空间，从而表示由稳定子本征值控制的条件位移。若稳定子本征值为第一取值，则目标模式保持在基准指针态；若稳定子本征值为第二取值，则目标模式被位移到可区分指针态。若具体实现中 $S_j=+1/-1$ 产生相反方向位移，则通过无条件补偿位移把两个结果等效转换为“基准指针态/可区分指针态”。支撑集合互不冲突且模式频率可分辨的多个 $S_j$ 可在同一脉冲窗口写入不同模式。通过边带选择、相干态相位判别或阈值判定，将两个指针态转化为综合征比特；在采用 Fock 态实现时，可进一步把二值指针态约束到 $n=0$ 与 $n=1$ 的占据窗口。

如图 3 所示为并行声子写入与逐模读出的执行流程。
\begin{figure}[h]
    \centering
    \includegraphics[width=0.75\textwidth]{Figure_3.png}
    \caption{并行声子写入与逐模读出流程图}
    \label{fig:flow}
\end{figure}

\subsubsection{3. 串扰抑制与错误标记}
控制模块根据模式频率、激光相位和离子参与向量设置脉冲包络，使非目标模式位移闭合，并在写入前确认目标模式初始化满足热占据阈值。若检测到模式热占据、边带漂移或写入概率异常，则将对应综合征标记为擦除并触发重测。

\subsubsection{4. 逐模读取综合征}
写入完成后，共享读出离子依次与每个目标模式耦合。通过红/蓝边带 $\pi$ 脉冲、旁带谱判别或量子逻辑读出把模式二值指针态映射为读出离子的内态，再以荧光测量得到综合征比特。读取顺序不改变数据离子状态；仅在相应模式的综合征已被读出并记录后，才对该模式执行再冷却、相干复位或下一轮初始化。

\subsubsection{5. 纠错反馈}
译码器接收由多个声子模式读出的综合征向量，执行 qLDPC 译码，并将 Pauli frame 更新发送至量子控制系统。

%==============================================================================
\section{本申请提案的关键点和欲保护点}
%==============================================================================

本申请提案主要包含以下关键创新点和保护点：

(1) \textbf{Tanner 图校验节点到声子本征模的映射}：将 qLDPC 稳定子校验节点直接映射为离子链可寻址正交运动模式，而非映射为每校验专用物理辅助离子。

(2) \textbf{声子二值指针态综合征存储}：以模式基准态和可区分位移态，或以 $n=0$ 和 $n=1$ 占据窗口，暂存稳定子奇偶信息，使声子模式成为临时综合征寄存器。

(3) \textbf{多稳定子并行写入}：利用模式正交性和频率选择性，通过多色状态依赖力同时执行多个高权重稳定子校验。

(4) \textbf{共享读出离子逐模读取}：使用一个或少量双种类读出离子对多个模式逐一进行边带读出，替代多个校验节点对应的专用辅助离子。

(5) \textbf{面向 qLDPC 的编译约束}：在编译阶段联合考虑 Tanner 图重叠关系、Pauli 基变换、模式频谱、参与向量（包括校验支撑集合内各离子在目标模式上的参与系数约束）、热占据、读出后模式复位和激光串扰，生成可执行的综合征提取脉冲表。

%==============================================================================
\section{与第三条中最接近的现有技术相比，本申请提案有何技术优点}
%==============================================================================

相比于现有以辅助离子和顺序门电路提取综合征的方案，本方案具有显著优势：

(1) \textbf{空间开销降低}：将 $K$ 个校验节点的临时存储由 $K$ 个专用物理辅助离子转化为 $K$ 个可寻址声子模式，并由一个或少量共享读出离子完成读取，减少离子数和寻址通道需求。

(2) \textbf{提取速度提升}：多个高权重稳定子可在同一脉冲窗口并行写入不同模式，降低 qLDPC 综合征提取深度。

(3) \textbf{适配高编码率 qLDPC}：本方案针对重叠校验和高权重稳定子设计，适合大逻辑容量离子阱量子处理器。

(4) \textbf{工程可实现性较强}：状态依赖力、边带读出、双种类离子和模式冷却均属于离子阱平台已有控制能力，本发明的创新在于资源映射和综合征存储方式。

%==============================================================================
\section{其他有助于理解本申请提案的技术资料}
%==============================================================================

(1) Panteleev, P. and Kalachev, G. ``Asymptotically Good Quantum and Locally Testable Classical LDPC Codes,'' 2021.
(2) Chen, W., et al. ``Quantum computation and simulation with vibrational modes of trapped ions,'' 2021.
(3) Hilder, J., et al. ``Fault-tolerant parity readout on a shuttling-based trapped-ion quantum computer,'' 2021.
(4) Ye, M. and Delfosse, N. ``Quantum error correction for long chains of trapped ions,'' 2025.

%==============================================================================
\section{本申请提案的侵权证据可获得性}
%==============================================================================

本提案的侵权行为具有较高的可检测性：
(1) \textbf{控制脉冲记录}：量子控制系统会保存多色边带频率、模式编号、脉冲包络、模式复位步骤和稳定子调度表，可反映校验节点到声子模式的映射。
(2) \textbf{读出时序证据}：若单个读出离子按多个运动模式逐一执行边带读出并输出 qLDPC 综合征，可作为实施本方案的重要线索。
(3) \textbf{硬件资源配置}：结合公开硬件配置、控制脉冲日志和读出时序，若系统在执行大量 qLDPC 校验时未配置对应数量的辅助离子，而依赖多个模式频率通道完成综合征暂存，可作为侵权判断线索。

\end{document}
